\section{Recipe 1: Xây dựng Hệ thống Phân loại Ảnh}
\label{sec:recipe_classification}

Phân loại ảnh là bài toán nền tảng nhất của thị giác máy tính, và cũng là điểm khởi đầu lý tưởng cho bất kỳ ai muốn xây dựng ứng dụng thực tế. Recipe này hướng dẫn bạn xây dựng một hệ thống phân biệt chó và mèo --- bài toán kinh điển của cộng đồng ML --- nhưng với một pipeline chuyên nghiệp: transfer learning từ ResNet-50, early stopping để tránh overfitting, và hàm suy luận có thể tích hợp ngay vào sản phẩm.

\subsection{Bước 1: Chuẩn bị cấu hình và dữ liệu}
\label{subsec:cls_data_prep}

Thay vì rải rác các hằng số khắp nơi trong code, hãy tập trung tất cả cấu hình vào một lớp \texttt{Config}. Thói quen này giúp thử nghiệm nhanh hơn và tránh lỗi khi thay đổi hyperparameter.

Dataset cần được tổ chức theo cấu trúc mà \texttt{ImageFolder} của PyTorch có thể nhận diện: mỗi thư mục con là một nhãn lớp, và tất cả ảnh của lớp đó nằm bên trong.

\begin{minted}[breaklines=true, fontsize=\small]{text}
data/dogs_vs_cats/
    cat/
        cat_001.jpg
        cat_002.jpg
        ...
    dog/
        dog_001.jpg
        dog_002.jpg
        ...
\end{minted}

\begin{minted}[breaklines=true, fontsize=\small]{python}
import torch
import torch.nn as nn
import torchvision.models as models
from torchvision import datasets, transforms
from torch.utils.data import DataLoader, random_split
import torch.optim as optim
from pathlib import Path
import copy

# Tập trung toàn bộ cấu hình vào một lớp
class Config:
    data_dir = "data/dogs_vs_cats"
    num_classes = 2
    batch_size = 32
    num_epochs = 20
    lr = 1e-3
    patience = 5  # early stopping: dừng nếu không cải thiện sau 5 epoch
    device = torch.device("cuda" if torch.cuda.is_available() else "cpu")

cfg = Config()
\end{minted}

\subsection{Bước 2: Data Augmentation}
\label{subsec:cls_augmentation}

Augmentation là kỹ thuật tăng cường dữ liệu bằng cách tạo ra các biến thể ngẫu nhiên của ảnh gốc trong quá trình huấn luyện. Tập train và tập validation cần các transform \textit{khác nhau}: tập train dùng augmentation ngẫu nhiên để tăng sự đa dạng, tập validation chỉ dùng resize và crop để đảm bảo đánh giá nhất quán.

\begin{minted}[breaklines=true, fontsize=\small]{python}
# Transform cho tập huấn luyện: augmentation ngẫu nhiên
train_transform = transforms.Compose([
    transforms.RandomResizedCrop(224),       # Cắt ngẫu nhiên và resize về 224x224
    transforms.RandomHorizontalFlip(),        # Lật ngang ngẫu nhiên (50%)
    transforms.ColorJitter(brightness=0.2, contrast=0.2),  # Thay đổi màu sắc
    transforms.ToTensor(),
    transforms.Normalize([0.485, 0.456, 0.406], [0.229, 0.224, 0.225]),  # Chuẩn hóa ImageNet
])

# Transform cho tập validation: cố định, không ngẫu nhiên
val_transform = transforms.Compose([
    transforms.Resize(256),
    transforms.CenterCrop(224),
    transforms.ToTensor(),
    transforms.Normalize([0.485, 0.456, 0.406], [0.229, 0.224, 0.225]),
])
\end{minted}

\begin{mdframed}[style=infobox]
\textbf{Tại sao dùng mean/std của ImageNet?}

Các giá trị \texttt{[0.485, 0.456, 0.406]} và \texttt{[0.229, 0.224, 0.225]} là mean và standard deviation của toàn bộ dataset ImageNet (hơn 1.2 triệu ảnh). Khi dùng một mô hình đã được pretrain trên ImageNet (như ResNet-50), bạn \textit{bắt buộc} phải chuẩn hóa input bằng chính các giá trị này --- vì đó là phân phối dữ liệu mà mô hình đã quen. Dùng giá trị khác sẽ làm các đặc trưng đã học mất ý nghĩa và giảm hiệu quả transfer learning đáng kể.
\end{mdframed}

\subsection{Bước 3: Nạp dữ liệu và phân chia tập train/val}
\label{subsec:cls_dataloader}

\begin{minted}[breaklines=true, fontsize=\small]{python}
# Nạp toàn bộ dataset với train_transform trước
full_dataset = datasets.ImageFolder(cfg.data_dir, transform=train_transform)

# Phân chia 80% train, 20% validation
train_size = int(0.8 * len(full_dataset))
val_size = len(full_dataset) - train_size
train_dataset, val_dataset = random_split(full_dataset, [train_size, val_size])

# Áp dụng val_transform cho tập validation
# Lưu ý: val_dataset.dataset trỏ đến full_dataset, cần ghi đè transform
val_dataset.dataset.transform = val_transform

train_loader = DataLoader(
    train_dataset, batch_size=cfg.batch_size, shuffle=True, num_workers=4
)
val_loader = DataLoader(
    val_dataset, batch_size=cfg.batch_size, shuffle=False, num_workers=4
)

print(f"Train: {len(train_dataset)} anh | Val: {len(val_dataset)} anh")
print(f"Cac lop: {full_dataset.classes}")
\end{minted}

\subsection{Bước 4: Xây dựng mô hình với Transfer Learning}
\label{subsec:cls_model}

Transfer learning là chiến lược sử dụng lại một mô hình đã được pretrain trên bộ dữ liệu lớn (ImageNet), sau đó tinh chỉnh (fine-tune) cho bài toán của mình. Đây là cách tiếp cận chuẩn mực trong thực tế vì tiết kiệm thời gian, giảm yêu cầu dữ liệu, và thường cho kết quả tốt hơn nhiều so với huấn luyện từ đầu.

Chiến lược ở đây là \textbf{đóng băng toàn bộ backbone} (các lớp tích chập) và chỉ huấn luyện \textbf{classification head} mới. Backbone đã có khả năng trích xuất đặc trưng hình ảnh xuất sắc; chúng ta chỉ cần dạy phần đầu phân loại biết cách dùng các đặc trưng đó cho bài toán mới.

\begin{minted}[breaklines=true, fontsize=\small]{python}
# Tải ResNet-50 với trọng số pretrained trên ImageNet
model = models.resnet50(weights=models.ResNet50_Weights.IMAGENET1K_V2)

# Đóng băng toàn bộ backbone: không cập nhật gradient
for param in model.parameters():
    param.requires_grad = False

# Thay thế classification head: in_features=2048 với ResNet-50
model.fc = nn.Sequential(
    nn.Dropout(0.3),
    nn.Linear(model.fc.in_features, cfg.num_classes)
)
model = model.to(cfg.device)

# Chỉ tối ưu các tham số của head mới (requires_grad=True)
optimizer = optim.AdamW(model.fc.parameters(), lr=cfg.lr, weight_decay=1e-4)
# CosineAnnealingLR: giảm lr theo đường cosine, giúp hội tụ mượt mà hơn
scheduler = optim.lr_scheduler.CosineAnnealingLR(optimizer, T_max=cfg.num_epochs)
criterion = nn.CrossEntropyLoss()

total_params = sum(p.numel() for p in model.parameters())
trainable_params = sum(p.numel() for p in model.parameters() if p.requires_grad)
print(f"Tong tham so: {total_params:,} | Tham so huan luyen: {trainable_params:,}")
\end{minted}

\subsection{Bước 5: Vòng lặp huấn luyện với Early Stopping}
\label{subsec:cls_training}

Early stopping là cơ chế dừng huấn luyện sớm khi mô hình ngừng cải thiện trên tập validation, giúp tránh overfitting và tiết kiệm thời gian. Chúng ta lưu lại bộ trọng số tốt nhất sau mỗi epoch cải thiện.

\begin{minted}[breaklines=true, fontsize=\small]{python}
best_acc = 0.0
best_model_wts = copy.deepcopy(model.state_dict())
patience_counter = 0

Path("models").mkdir(exist_ok=True)

for epoch in range(cfg.num_epochs):
    # --- Pha huấn luyện ---
    model.train()
    train_loss = 0.0
    for images, labels in train_loader:
        images, labels = images.to(cfg.device), labels.to(cfg.device)
        optimizer.zero_grad()
        outputs = model(images)
        loss = criterion(outputs, labels)
        loss.backward()
        optimizer.step()
        train_loss += loss.item()

    # --- Pha đánh giá ---
    model.eval()
    correct, total = 0, 0
    with torch.no_grad():
        for images, labels in val_loader:
            images, labels = images.to(cfg.device), labels.to(cfg.device)
            outputs = model(images)
            _, predicted = outputs.max(1)
            total += labels.size(0)
            correct += predicted.eq(labels).sum().item()

    val_acc = correct / total
    scheduler.step()

    print(
        f"Epoch {epoch+1}/{cfg.num_epochs} | "
        f"Train Loss: {train_loss/len(train_loader):.4f} | "
        f"Val Acc: {val_acc:.4f}"
    )

    # Lưu mô hình tốt nhất và kiểm tra early stopping
    if val_acc > best_acc:
        best_acc = val_acc
        best_model_wts = copy.deepcopy(model.state_dict())
        torch.save(model.state_dict(), "models/best_classifier.pth")
        patience_counter = 0
        print(f"  -> Mo hinh tot nhat duoc luu (acc={best_acc:.4f})")
    else:
        patience_counter += 1
        if patience_counter >= cfg.patience:
            print(f"Early stopping tai epoch {epoch+1}")
            break

# Nạp lại trọng số tốt nhất
model.load_state_dict(best_model_wts)
print(f"\nHuan luyen hoan tat. Val Acc tot nhat: {best_acc:.4f}")
\end{minted}

\subsection{Bước 6: Hàm suy luận}
\label{subsec:cls_inference}

Sau khi huấn luyện, bạn cần một hàm suy luận gọn gàng, có thể tái sử dụng ở bất kỳ đâu trong ứng dụng.

\begin{minted}[breaklines=true, fontsize=\small]{python}
def predict_image(image_path, model, transform, class_names):
    """Dự đoán nhãn lớp và độ tự tin cho một ảnh."""
    from PIL import Image
    model.eval()
    image = Image.open(image_path).convert("RGB")
    tensor = transform(image).unsqueeze(0).to(cfg.device)  # Thêm batch dimension
    with torch.no_grad():
        output = model(tensor)
        prob = torch.softmax(output, dim=1)[0]  # Chuyển logits sang xác suất
        pred_idx = prob.argmax().item()
    return class_names[pred_idx], float(prob[pred_idx])

# Sử dụng
class_names = full_dataset.classes  # ['cat', 'dog']
pred_class, confidence = predict_image(
    "test.jpg", model, val_transform, class_names
)
print(f"Du doan: {pred_class} (do tu tin: {confidence:.2%})")
\end{minted}

\begin{mdframed}[style=infobox]
\textbf{Nâng cao độ chính xác: ba kỹ thuật nên thử tiếp theo}

Sau khi recipe cơ bản hoạt động, hãy thử các cải tiến sau theo thứ tự:

\begin{enumerate}
    \item \textbf{Mở đóng băng backbone dần dần (gradual unfreezing):} Sau khi head đã ổn định (khoảng 5 epoch), mở đóng băng backbone từ các lớp cuối lên đầu với learning rate rất nhỏ (1e-5 đến 1e-4). Kỹ thuật này thường cải thiện thêm 1--3\% accuracy.

    \item \textbf{Mixup augmentation:} Trộn hai ảnh và nhãn của chúng với tỉ lệ ngẫu nhiên $\lambda \sim \text{Beta}(0.2, 0.2)$: $\tilde{x} = \lambda x_i + (1-\lambda) x_j$. Buộc mô hình học các biểu diễn mượt mà hơn và giảm overfitting đáng kể.

    \item \textbf{Label smoothing:} Thay vì nhãn one-hot cứng nhắc $[0, 1]$, dùng $[0.05, 0.95]$. Ngăn mô hình trở nên quá tự tin và cải thiện calibration. Chỉ cần thay \texttt{CrossEntropyLoss()} bằng \texttt{CrossEntropyLoss(label\_smoothing=0.1)}.
\end{enumerate}
\end{mdframed}
